\chapter{Tecnologías utilizadas y criterio de elección}
\label{cap:tecnologias}

Este capítulo aterriza la solución en tecnologías concretas. Su objetivo principal es describir qué es cada tecnología y qué papel desempeña dentro del sistema completo. Cuando resulta útil, se menciona también por qué encaja en el proyecto, pero el detalle de cómo se ha combinado con otras piezas se deja para el Capítulo~\ref{cap:implementacion}. También resulta importante distinguir entre herramientas preexistentes reutilizadas en el proyecto y componentes desarrollados específicamente en este TFG.

% ============================================================
\section{Python como lenguaje de integración principal}
\label{sec:python}
% ============================================================

El lenguaje vertebrador del proyecto es Python. La razón principal no es solo su popularidad, sino su capacidad para actuar como lenguaje de integración entre bibliotecas de IA, servicios web, automatización y procesamiento documental. En este TFG, Python permite conectar en un mismo ecosistema el backend, el indexador, el pipeline de OpenWebUI, la integración con SharePoint, el scraping, el OCR y el servidor MCP.

Conviene separar aquí dos planos. Por un lado, hay paquetes que se usan directamente en el código desarrollado en el TFG, como FastAPI, qdrant-client, sentence-transformers o msal. Por otro, hay dependencias que vienen dadas por herramientas ya reutilizadas en el stack, como parte del ecosistema de OpenWebUI, Docker o Grafana. Esta distinción es importante porque el proyecto no consiste en construir esas plataformas de base, sino en apoyarse en ellas y extenderlas donde era necesario.

Además, buena parte de las herramientas reutilizadas en el proyecto exponen su mejor soporte precisamente en este lenguaje. Entre los paquetes y bibliotecas más importantes utilizados se encuentran:

\begin{itemize}
    \item \textbf{FastAPI} y \textbf{Uvicorn} para los servicios web y la API REST.
    \item \textbf{Pydantic} para validación de datos y configuración tipada.
    \item \textbf{sentence-transformers} para generación de embeddings densos y \textbf{fastembed} para los vectores dispersos (BM25).
    \item \textbf{qdrant-client} para interacción con la base de datos vectorial.
    \item \textbf{msal} y clientes asociados para autenticación y acceso a Microsoft Graph.
    \item \textbf{playwright} para automatización de navegador.
    \item \textbf{trafilatura}, \textbf{readability} y \textbf{BeautifulSoup} para extracción de contenido web.
    \item \textbf{PaddleOCR} para OCR de documentos escaneados y \textbf{Ray} para paralelizar su procesamiento.
    \item \textbf{duckduckgo\_search} para búsqueda web controlada.
    \item \textbf{fastmcp} para el servidor MCP del BOE.
    \item \textbf{prometheus-client} para exponer métricas del backend.
\end{itemize}

Esta elección también ayuda a explicar la naturaleza del trabajo realizado: el proyecto no consiste en inventar estas bibliotecas, sino en utilizarlas de forma coordinada y construir sobre ellas la lógica que faltaba para el caso de uso corporativo.

% ============================================================
\section{OpenWebUI: interfaz conversacional reutilizada}
\label{sec:openwebui}
% ============================================================

OpenWebUI\footnote{\url{https://openwebui.com/}} es la interfaz conversacional utilizada como punto de entrada del sistema \cite{openwebui2024docs}. Es una herramienta externa, ya existente, que he reutilizado para no invertir el esfuerzo del TFG en construir una interfaz de chat desde cero. Esta decisión permite concentrar el trabajo en la parte realmente diferencial del proyecto: la integración con las fuentes documentales y la lógica del asistente.

\subsection{Arquitectura interna de OpenWebUI}

OpenWebUI se estructura como una aplicación web con tres bloques principales:

\begin{itemize}
    \item \textit{Frontend}: interfaz tipo chat accesible desde navegador, con historial de conversaciones, carga de adjuntos y streaming de respuestas.
    \item \textit{Backend}: lógica de autenticación, persistencia de chats y gestión de modelos.
    \item \textit{Base de datos}: almacenamiento de usuarios, conversaciones y metadatos.
\end{itemize}

Para el lector, lo importante es entender su papel dentro del proyecto: OpenWebUI aporta la experiencia de usuario, pero no constituye la inteligencia principal del sistema. Esa inteligencia se añade mediante el pipeline desarrollado en este TFG. En consecuencia, cuando en esta memoria se habla de ``integrar autenticación corporativa con la interfaz conversacional'', no se está describiendo la creación de una interfaz nueva, sino la reutilización de OpenWebUI junto con su conexión al proveedor corporativo de identidad y su ampliación mediante un pipeline propio.

\subsection{Por qué OpenWebUI}

La elección de OpenWebUI se justifica porque ya ofrece varias capacidades deseables en un entorno corporativo: interfaz madura, persistencia de conversaciones, compatibilidad con modelos locales y soporte de autenticación. Frente a hacer una interfaz propia, reduce tiempo de desarrollo y riesgo técnico.

\begin{table}[H]
\centering
\caption{Comparativa de interfaces de usuario para LLMs locales.}
\label{tab:openwebui-comparativa}
\begin{tabular}{lcccc}
\toprule
\textbf{Criterio} & \textbf{OpenWebUI} & \textbf{Gradio} & \textbf{Streamlit} & \textbf{ChatUI} \\
\midrule
UI tipo ChatGPT              & \checkmark & Parcial   & Parcial  & \checkmark \\
Compatibilidad con Ollama    & Nativa     & Manual    & Manual   & Nativa \\
Sistema de Pipelines/Plugins & \checkmark & $\times$  & $\times$ & $\times$ \\
Gestión de usuarios          & \checkmark & $\times$  & $\times$ & \checkmark \\
Persistencia de chats        & \checkmark & $\times$  & $\times$ & \checkmark \\
SSO / OIDC                   & \checkmark & $\times$  & $\times$ & $\times$ \\
Código abierto               & \checkmark & \checkmark & \checkmark & \checkmark \\
Carga de archivos adjuntos   & \checkmark & \checkmark & \checkmark & Parcial \\
\bottomrule
\end{tabular}
\end{table}

El factor decisivo es el sistema de Pipelines, porque permite insertar lógica personalizada sin modificar el núcleo de la herramienta. Gracias a ello, la interfaz sigue siendo una pieza reutilizada, pero el comportamiento final del asistente pasa a estar controlado por desarrollo propio. El funcionamiento de este mecanismo se describe en el Capítulo~\ref{cap:implementacion} y su implementación detallada en el Anexo~\ref{anexo:pipeline}.

% ============================================================
\section{Ollama: ejecución local de modelos}
\label{sec:ollama}
% ============================================================

Ollama\footnote{\url{https://ollama.com/}} es la herramienta externa utilizada para ejecutar modelos de lenguaje localmente. Su función es cargar modelos, servirlos por API y gestionar su uso en el hardware disponible \cite{ollama2024docs}. Se ha elegido porque facilita el despliegue de modelos locales sin construir una capa de inferencia propia, gestiona la VRAM de manera operativa (carga y descarga modelos según necesidad), expone una API sencilla y encaja bien como un servicio más dentro de Docker Compose.

\subsection{Modelos desplegados}

En el sistema se emplean varias familias de modelos ya existentes, cada una con una función concreta:

\begin{table}[H]
\centering
\caption{Modelos de IA desplegados en Ollama y su función en el sistema.}
\label{tab:modelos-ollama}
\resizebox{\textwidth}{!}{%
\begin{tabular}{llcc}
\toprule
\textbf{Modelo} & \textbf{Función} & \textbf{Cuantización} & \textbf{VRAM} \\
\midrule
rag-qwen-ft:latest    & Modelo principal del alias JARVIS para chat y RAG & LoRA FT & $\sim$8 GB \\
qwen2.5:32b-instruct-q4\_K\_M & Fallback de texto y comparación base & Q4\_K\_M & $\sim$19 GB \\
qwen2.5-coder:32b     & Modelo auxiliar instalado para tareas de código & --- & $\sim$19 GB \\
qwen2.5vl:7b          & Análisis de imágenes y OCR multimodal & --- & $\sim$6 GB \\
\bottomrule
\end{tabular}%
}
\end{table}

La tabla refuerza una idea importante: los modelos no se han creado en este TFG, pero sí se han seleccionado, desplegado, probado y asignado a tareas concretas dentro de una arquitectura coherente. En el despliegue actual del proyecto, la familia Qwen 2.5 concentra las funciones principales de texto, visión y adaptación al dominio.

\subsection{Stack de modelos en producción}
\label{subsec:stack-modelos}

La elección del LLM principal no fue inmediata: evolucionó a lo largo del proyecto en cuatro etapas ---de \texttt{llama3.1:8b} en cuantización Q4 y Q8, a Qwen 2.5 7B base y, finalmente, al modelo ajustado con LoRA---, cada una motivada por limitaciones reales de calidad, ventana de contexto o rendimiento en español detectadas durante el uso. Esa trayectoria completa, con sus causas técnicas (incluido el problema de los \textit{stop tokens} de Qwen), se recoge en el Anexo~\ref{anexo:evolucion}. La Tabla~\ref{tab:stack-llm-produccion} resume el stack de modelos resultante, que combina el LLM fine-tuned, un modelo grande para análisis, un modelo de visión y los modelos auxiliares de recuperación (embeddings densos, dispersos y reranker).

\begin{table}[H]
\centering
\caption{Stack de modelos en producción (versión TFG y sistema empresarial).}
\label{tab:stack-llm-produccion}
\resizebox{\textwidth}{!}{%
\begin{tabular}{lllcc}
\toprule
\textbf{Alias / Rol} & \textbf{Modelo} & \textbf{Función} & \textbf{VRAM} & \textbf{Contexto} \\
\midrule
JARVIS / europav-IA & \texttt{rag-qwen-ft:latest}            & Chat RAG principal con fine-tuning LoRA & $\sim$8\,GB  & 32\,K \\
qwen2.5-32b         & \texttt{qwen2.5:32b-instruct-q4\_K\_M} & Fallback y análisis complejos         & $\sim$19\,GB & 32\,K \\
qwen2.5vl           & \texttt{qwen2.5vl:7b}                  & Visión e interpretación de imágenes & $\sim$6\,GB  & 32\,K \\
Embeddings (densos) & \texttt{paraphrase-multilingual-MiniLM-L12-v2} & Vectorización semántica de texto (CPU) & ---         & 512 tok. \\
Sparse (BM25)       & \texttt{Qdrant/bm25} (\texttt{fastembed}) & Vectores dispersos para búsqueda léxica (CPU) & --- & --- \\
Reranker            & \texttt{cross-encoder/ms-marco-MiniLM-L-6-v2} & Reordenación de candidatos por relevancia (CPU) & --- & --- \\
\bottomrule
\end{tabular}%
}
\end{table}

El stack de producción cubre la generación con modelos especializados e independientes para no saturar la VRAM: el modelo fine-tuned para RAG corporativo, un modelo de 32\,B para análisis complejo y un modelo multimodal para visión. Tanto el modelo fine-tuned (JARVIS) como el de 32\,B heredan la ventana de contexto de 32\,K tokens de la familia Qwen 2.5. La recuperación se apoya en tres modelos ligeros que se ejecutan en CPU: el de embeddings densos, el de vectores dispersos BM25 y el reranker Cross-Encoder. LiteLLM actúa como proxy unificado con enrutamiento automático, caché Redis y fallback entre modelos, y mantiene aliases de compatibilidad hacia atrás (\texttt{llama3.1-8b}~$\rightarrow$~\texttt{qwen2.5-32b}) para integraciones anteriores.

\subsection{Query Processor: enrutamiento inteligente y expansión de consultas}
\label{subsec:query-processor}

El \texttt{QueryProcessor} (\texttt{core/query\_processor.py}) es un componente transversal que intercepta cada consulta del usuario antes de que llegue a la base de datos vectorial, enriqueciéndola en tres dimensiones:

\paragraph{Detección de intención (\textit{Intent Detection}).}
Clasifica la consulta en cuatro tipos mediante expresiones regulares sobre el texto en español:
\begin{itemize}
    \item \textbf{FACTUAL} --- pregunta puntual (``¿qué es...?'', ``¿cuándo...?''): \texttt{top\_k=5}, búsqueda densa.
    \item \textbf{PROCEDURAL} --- cómo hacer algo (``¿cómo...?'', ``pasos para...''): \texttt{top\_k=8}, búsqueda híbrida.
    \item \textbf{ANALYTICAL} --- análisis o comparación (``diferencia entre...'', ``resume...''): \texttt{top\_k=12}, búsqueda híbrida.
    \item \textbf{CONVERSATIONAL} --- saludo o chat general: \texttt{top\_k=0}, sin búsqueda RAG.
\end{itemize}

\paragraph{Enrutamiento de modelos (\textit{Smart Model Routing}).}
La intención detectada determina qué modelo LLM genera la respuesta final: las consultas FACTUAL y PROCEDURAL van a \textbf{JARVIS} (\texttt{rag-qwen-ft:latest}: fine-tuned, rápido, citas correctas), y las ANALYTICAL a \textbf{qwen2.5-32b} (ventana de 32\,K, mayor capacidad de razonamiento y síntesis multi-documento). Es el principio de \textit{right tool for the job}.

\paragraph{Expansión de consultas (\textit{Query Expansion}).}
Para consultas no conversacionales, el \texttt{QueryProcessor} genera dos reformulaciones adicionales usando JARVIS como LLM auxiliar (con timeout de 12\,s). Las tres variantes se lanzan en paralelo contra Qdrant y sus resultados se fusionan eliminando duplicados antes de pasarlos al reranker, lo que mejora el \textit{recall} semántico cuando el usuario usa terminología distinta a la del documento:

\begin{center}
\textit{``normas seguridad escaleras''} $\xrightarrow{\text{expansión}}$
\textit{``requisitos técnicos escaleras portátiles''},
\textit{``normativa PRL trabajo en altura escaleras''}
\end{center}

La combinación de los tres mecanismos ---intent detection, model routing y query expansion--- convierte el pipeline RAG en un sistema adaptativo que ajusta automáticamente su estrategia de recuperación y generación según la naturaleza de cada pregunta, sin intervención del usuario.

% ============================================================
\section{Qdrant: base de datos vectorial}
\label{sec:qdrant}
% ============================================================

Qdrant\footnote{\url{https://qdrant.tech/}} es la base de datos vectorial empleada como almacén principal del conocimiento documental del sistema \cite{qdrant2024docs}. Es una herramienta externa, pero su estructura, colecciones y criterios de uso se definen dentro del proyecto. De sus capacidades, el sistema aprovecha el \textit{índice HNSW} (búsquedas vectoriales rápidas sobre grandes volúmenes de fragmentos), la \textit{búsqueda híbrida} (combina similitud semántica densa y coincidencia léxica dispersa), el \textit{filtrado por metadatos} (segmentación por fuente, documento o departamento) y el soporte \textit{multi-colección} (separación de documentos por contexto de uso). Para el proyecto, esto significa que Qdrant no es solo un repositorio técnico, sino el componente que hace posible que el sistema encuentre evidencia relevante antes de responder.

% ============================================================
\section{FastAPI: backend del sistema}
\label{sec:fastapi}
% ============================================================

FastAPI\footnote{\url{https://fastapi.tiangolo.com/}} es el framework seleccionado para construir los servicios backend en Python \cite{fastapi2024docs}. A diferencia de OpenWebUI, Ollama o Qdrant, que son herramientas reutilizadas, aquí sí aparece una parte importante de desarrollo propio: los servicios implementados con FastAPI contienen gran parte de la lógica central del sistema. Se ha elegido por su facilidad para construir APIs tipadas y mantenibles, su soporte asíncrono (útil cuando el sistema depende de múltiples llamadas a servicios externos o internos) y su integración natural con el ecosistema Python (IA, scraping, OCR y autenticación en el mismo entorno). En este TFG, FastAPI sustenta el backend RAG, el indexador de documentos y parte de las integraciones auxiliares.

% ============================================================
\section{LiteLLM: proxy unificado de modelos}
\label{sec:litellm}
% ============================================================

LiteLLM\footnote{\url{https://docs.litellm.ai/}} es una herramienta externa que se utiliza como capa intermedia entre los servicios propios y los modelos servidos por Ollama \cite{litellm2024docs}. Su utilidad está en unificar llamadas, aplicar alias de modelos, gestionar caché y simplificar el encaminamiento técnico. Desde la perspectiva del lector no técnico, su papel puede resumirse así: en lugar de que cada parte del sistema tenga que conocer todos los detalles de cada modelo, LiteLLM ofrece un punto único y homogéneo de acceso.

% ============================================================
\section{Docker y orquestación de contenedores}
\label{sec:docker}
% ============================================================

Docker\footnote{\url{https://www.docker.com/}} y Docker Compose son la base del despliegue reproducible \cite{dockerCompose2026docs}. Se utilizan para empaquetar todos los componentes y ponerlos en marcha como un conjunto coordinado. Sus ventajas en este proyecto son la \textit{reproducibilidad} (el mismo sistema se despliega de forma controlada en distintos entornos), el \textit{aislamiento} (cada servicio mantiene sus dependencias separadas), la \textit{seguridad perimetral} (solo ciertos servicios se exponen al exterior) y el \textit{acceso a GPU bajo demanda} (se reserva aceleración solo para los servicios que la necesitan). En el contexto del TFG, esto también es parte de la contribución: no solo se han elegido tecnologías, sino que se han dejado integradas dentro de una arquitectura desplegable de forma repetible.

% ============================================================
\section{Herramientas de procesamiento documental y web}
\label{sec:herramientas-procesamiento}
% ============================================================

Esta sección reúne herramientas ya existentes que no han sido desarrolladas en el proyecto, pero que sí se han integrado y combinado para cubrir necesidades concretas del sistema.

\subsection{PaddleOCR: reconocimiento óptico de caracteres}

PaddleOCR\footnote{\url{https://github.com/PaddlePaddle/PaddleOCR}} es un motor OCR externo utilizado para extraer texto de PDFs escaneados \cite{paddleocr2024docs}. No se ha implementado un OCR propio; el trabajo realizado consiste en decidir cuándo usarlo, cómo integrarlo en la ingesta y cómo hacerlo convivir con el resto de componentes. La paralelización de su procesamiento se apoya en \textbf{Ray}.

\subsection{Playwright: automatización de navegadores}

Playwright\footnote{\url{https://playwright.dev/}} es una herramienta externa de automatización de navegador \cite{playwright2024docs}. Se utiliza porque muchas páginas web modernas dependen de JavaScript y no pueden procesarse con una simple petición HTTP. Su papel dentro del proyecto es permitir obtener el contenido real que ve el usuario en páginas dinámicas.

\subsection{Trafilatura: extracción de contenido web}

Trafilatura \cite{barbaresi2021trafilatura} es una biblioteca externa especializada en extraer el contenido principal de una página web, eliminando menús, barras laterales y ruido. En este TFG no se ha desarrollado un extractor propio desde cero; se ha integrado Trafilatura como primera opción dentro de una cadena de extracción con \textit{fallbacks}, porque resuelve bien el problema de aislar el contenido relevante.

\subsection{SentenceTransformers y reranker: recuperación}

SentenceTransformers \cite{reimers2019sentence,sbertDocs} es la biblioteca utilizada para generar los embeddings densos, con el modelo \texttt{paraphrase-multilingual-MiniLM-L12-v2}, elegido por su equilibrio entre tamaño, soporte multilingüe y facilidad de despliegue. La recuperación se completa con dos piezas adicionales: los vectores dispersos \textbf{BM25} (\texttt{Qdrant/bm25} vía \texttt{fastembed}), que aportan coincidencia léxica exacta, y un \textbf{reranker} de tipo \textit{Cross-Encoder} (\texttt{cross-encoder/ms-marco-MiniLM-L-6-v2}) que reordena los candidatos por relevancia. Aquí también es importante la distinción metodológica: los modelos y librerías ya existían; lo que aporta el proyecto es su integración en el pipeline RAG y su ajuste al caso de uso documental.

% ============================================================
\section{Stack de monitorización}
\label{sec:stack-monitorizacion}
% ============================================================

Para operar un sistema de este tipo no basta con que funcione: también hay que entender cómo se comporta. Por ello se incorpora un stack de observabilidad formado por herramientas ya existentes, configuradas y conectadas dentro del proyecto.

\subsection{Prometheus}

Prometheus\footnote{\url{https://prometheus.io/}} recoge métricas temporales del backend y de otros servicios \cite{prometheus2026docs}. Su utilidad en este TFG es permitir medir latencia, tasa de errores y patrones de uso en lugar de basar las decisiones de optimización en impresiones subjetivas.

\subsection{Grafana}

Grafana\footnote{\url{https://grafana.com/}} visualiza las métricas recogidas por Prometheus mediante paneles comprensibles \cite{grafana2026docs}. En el proyecto se ha utilizado para construir paneles específicos de backend, GPU y bases de datos. Esto no debe entenderse como el desarrollo de una herramienta nueva, sino como el diseño de la capa de observación necesaria para operar el sistema con criterio.

\subsection{NVIDIA DCGM Exporter}

NVIDIA DCGM Exporter\footnote{\url{https://github.com/NVIDIA/dcgm-exporter}} expone métricas de la GPU en formato compatible con Prometheus \cite{dcgmExporter2026docs}. Su presencia se justifica porque en un sistema que ejecuta modelos y OCR en local, la GPU es un recurso crítico y debe monitorizarse con el mismo nivel de detalle que el backend.

\subsection{Loki y Promtail}

Mientras Prometheus se ocupa de las métricas numéricas, \textbf{Loki} y \textbf{Promtail} cubren los \textit{logs}: Promtail recolecta los registros de todos los contenedores y los envía a Loki, que los indexa y los hace consultables desde el mismo Grafana. Esto permite correlacionar un pico de latencia con las líneas de log exactas de ese instante, unificando métricas y logs en una sola herramienta de diagnóstico.

% ============================================================
\section{Otras tecnologías del ecosistema}
\label{sec:otras-tecnologias}
% ============================================================

\begin{itemize}
    \item \textit{PostgreSQL}: almacena usuarios, conversaciones y metadatos de OpenWebUI.
    \item \textit{Redis}: se emplea como caché de respuestas a través de LiteLLM.
    \item \textit{MinIO}: conserva documentos binarios como respaldo y fuente de reindexación.
    \item \textit{NGINX}: actúa como proxy inverso y punto único de entrada.
    \item \textit{MSAL} (\textit{Microsoft Authentication Library}): permite autenticación con Azure AD y acceso autenticado a SharePoint mediante Microsoft Graph.
    \item \textit{Ray}: reparte el procesamiento OCR entre varios \textit{workers} durante la ingesta.
\end{itemize}

El diagrama general que resume cómo se conectan todas estas tecnologías ---el flujo de consulta, la lógica propia del TFG, las integraciones externas y la capa de observabilidad--- se presenta en el Anexo~\ref{anexo:evolucion} (Figura~\ref{fig:flujo-tecnologias}), por su tamaño apaisado.
